\section{Detailed Primary Comparisons}\label{app:primary-context}
The following table and plot retain the full six-configuration primary
comparison and all six painter pairs discussed in the main text. The target
and uncertainty conventions are unchanged.

\begin{table}[tbp]
\caption{Agreement with the full-frame reference target. $\beta=1$ matches its aligned amplitude; larger is not necessarily better. Expected $D$ is 0 for exact agreement and 1 for no painter distinctions. Positive $\Delta D$ favors FLUX.2 Max. Table brackets give approximate 95\% simultaneous intervals from the prespecified family of six aligned responses and 15 scene-paired model contrasts. FLUX's lowest $D$ estimate has an unadjusted 95\% interval $[0.586,1.016]$, spanning the no-distinction value 1; this interval is descriptive.}
\label{tab:specificity-models}

\centering\small\setlength{\tabcolsep}{4pt}
\begin{tabularx}{\linewidth}{@{}>{\raggedright\arraybackslash}p{.24\linewidth}>{\raggedright\arraybackslash}Xr>{\raggedright\arraybackslash}X@{}}
\toprule
Model & Aligned response $\beta$ & \shortstack{Uncalibrated\\error $D$} & $\Delta D$ versus FLUX\\
\midrule
GPT Image 1 & $0.938\;[0.709,1.168]$ & $1.572$ & $0.771\;[-0.037,1.580]$\\[3pt]
GPT Image 2 & $0.999\;[0.893,1.104]$ & $1.226$ & $0.425\;[-0.161,1.011]$\\[3pt]
GPT Image 2.5 Flare & $0.772\;[0.680,0.864]$ & $1.738$ & $0.937\;[0.256,1.618]$\\[3pt]
GPT Image 2.5 Sunburst & $0.824\;[0.680,0.968]$ & $1.641$ & $0.840\;[0.058,1.623]$\\[3pt]
Nano Banana 2 & $0.442\;[0.256,0.629]$ & $1.109$ & $0.309\;[-0.847,1.464]$\\[3pt]
FLUX.2 Max & $0.470\;[0.239,0.701]$ & $0.801$ & $0$ (baseline)\\[3pt]
\bottomrule\end{tabularx}

\end{table}

\begin{figure}[!htbp]
\centering\includegraphics[width=\linewidth]{specificity_artist_pairs.pdf}
\caption{Descriptive painter-pair estimates (no pairwise tests). M, S, P and C denote Monet, Sisley,
Pissarro and C\'ezanne. Pair aligned response 0 means no aligned component,
1 matches reference strength, values above 1 exaggerate it, and negative values
indicate the opposite direction.
Expected error is 0 for exact agreement and 1 for no distinction. Each pair uses
all 31 features and is normalized by its squared reference distance.
Negative error can arise from repeat correction,
not better-than-exact agreement.}
\label{fig:artist-pairs}
\end{figure}
\FloatBarrier

\subsection{Detailed interpretation of the original comparisons}
\paragraph{A Response in the Right Direction Can Still Miss the Target}
All 1,008 outputs yield valid measurements. With the full-frame
reference target, every model's aligned-response interval lies above zero
(Appendix Table~\ref{tab:specificity-models}). A purely common additive response is
therefore insufficient to explain the named conditions under the stated
inference model. Yet positive alignment does not ensure low error. GPT Image 2
nearly matches the reference-aligned strength, $\beta=.999\;[.893,1.104]$,
while its total error is $D=1.226$: it adds substantial differences beyond
that aligned component.




FLUX.2 Max has the lowest error estimate, $.801$, despite a smaller aligned
response, $.470$. Adjusted comparisons resolve higher error for Flare and Sunburst
than for FLUX; the other 13 intervals include zero (Appendix~\ref{app:results},
Table~\ref{tab:all-specificity-pairs}).
These use full-frame references;
the Sunburst result changes after source correction
(Section~\ref{sec:source-correction}). No model is established to beat the
no-contrast benchmark of one. Thus the study detects a
reference-aligned response more clearly than it establishes low total mismatch.

Appendix~\ref{app:geometry} visualizes the same labeled contrasts in two
reference-fitted axes; Appendix Figure~\ref{fig:artist-pairs} uses all 31 features.



\paragraph{Uneven Painter Distinctions in the 31 Features}
Retrospective pair diagnostics extend the lightness example in
Section~\ref{sec:contrast-intuition} to all 31 features (Appendix Figure~\ref{fig:artist-pairs}).
Monet--Sisley aligned responses range from $-.249$ to $.129$ in five models,
versus $.402$ for GPT Image 2. Yet GPT Image 1 has lower estimated error for this pair
($.56$ versus $1.16$): stronger alignment does not ensure lower mismatch.
Omitting C\'ezanne lowers FLUX's overall response from $.470$ to $.096$ and
GPT Image 2's from $.999$ to $.651$.

Swapping Monet and Sisley lowers error for Flare, Sunburst and Nano Banana 2.
Reversing a negatively aligned pair difference lowers error without changing its size.




\paragraph{Which Comparisons Survive Reference Changes?}
Within the 31-feature metric, FLUX retains the lowest uncalibrated scene-error point
estimate across scene deletions, measurement windows, reference targets and
feature-weighting sensitivities (Appendices~\ref{app:declared-sensitivities}
and~\ref{app:diagnostics}).

Source handling has a more consequential effect on statistical comparisons.
AI assistants inspected all reference and development reproductions, selecting
crops that exclude peripheral artifacts in 90 reference and 41 development images.
The audit used no human verification. After cropping and refitting development scaling,
FLUX's error estimate changes from $.801$ to $.872$. Corrected aligned-response
intervals remain above zero for all six models, and GPT Image 2 retains the lowest
calibrated estimate.

Appendix Table~\ref{tab:source-comparison} makes the changed comparisons explicit. With
corrected regions and refitted scaling, the Sunburst interval crosses zero,
Flare still has higher error than FLUX and the GPT Image 1 interval moves above zero. The latter
is a retrospective sensitivity, not a new confirmatory finding. Source correction
preserves FLUX's lowest point estimate while changing the evidence for particular
model differences.
Weak Monet--Sisley alignment also persists after correction, but individual
signs are unstable: Flare changes from $-.011$ to $.019$, GPT Image 1 from
$.074$ to $-.053$, and FLUX from $.129$ to $.019$. The label-swap observation
above therefore concerns the original reference target.






A separate check uses class-specific reference means for 11 water, built and
land briefs, excluding three mixed briefs. Replacing title-derived reference
labels with visual labels retains FLUX's lowest error estimate.
Appendix~\ref{app:diagnostics} gives the controls and audit.\par



